\headline{{\textbf{\Large\sffamily Bonsai Almanac:}}\\
          {\textbf{\large\sffamily A Monographic Guide to Itoigawa Juniper}}}
\label{2025-07-itoigawa}
\medskip
\topic{A Shohin Plan and a Bigger Pot}
\noindent
Some trees sneak up on you.
For me, that tree was an \textit{Itoigawa juniper}.
I picked it up at last year's Crawley Bonsai Fest, from Herons.
It came in a tiny plastic pot and, although the foliage was healthy, I had the feeling something was off.
I noticed water pooling on the surface and realised it wasn't saturating the rootball properly.

Not wanting to risk a root prune on a juniper, I opted for a slip\--potting instead.
I moved it into a 3D\--printed training pot (146\texttimes146\texttimes70mm) with better drainage and drilled holes through the existing soil to help water reach deeper.
Still, it didn't make a significant difference.
So I've printed a slightly larger 190\texttimes190\texttimes70mm pot, gently teased the sides of the root ball and filled it up with draining soil.
With no pruning and more space, I'm hoping for some lower trunk thickening before I begin its development into a compact shohin.

This journey led me to learn more about this fascinating species.
What follows is a comprehensive guide to cultivating \textit{Juniperus chinensis `Itoigawa'} as a bonsai.
I hope it helps anyone considering working with this vigorous and rewarding tree.

\medskip
\topic{\large Cultivating Juniperus chinensis `Itoigawa' as a Bonsai: A Complete Guide}
\noindent
\textit{Juniperus chinensis `Itoigawa'} is one of the most prized cultivars in the bonsai world, especially among lovers of shohin and compact trees.
With its tight, vivid green foliage and elegant, naturally compact growth habit, it lends itself beautifully to high\--detail refinement.
Originally cultivated in the Itoigawa region of Japan, it's a selected variety of the Chinese juniper that offers both aesthetic and practical advantages.

This article explores its botanical features, ideal cultivation conditions, propagation methods, styling techniques, and more---including a detailed training timeline and common mistakes to avoid.

\begin{window}[%
    2,l,%
    \includegraphics[width=1.0\columnwidth]{2025_07/res/event/my_itoigawa},%
    \centerline{\emph{The Itoigawa I bought at Crawley Bonsai Show.}}%
]
\end{window}
\columnbreak

\medskip
\topic{Botanical Features and Visual Qualities}
\noindent
\textit{Itoigawa juniper} belongs to the Cupressaceae family and is a sub\--variety of \textit{Juniperus chinensis}.
This cultivar is recognised for:

\begin{itemize}
    \item \textbf{Foliage:} Dense, scale\--like, emerald\--green foliage that resists reversion to needle growth even under stress.

    \item \textbf{Bark:} Attractive reddish\--brown bark that exfoliates in ribbons, revealing pale streaks underneath---ideal for deadwood work.

    \item \textbf{Growth Habit:} Naturally compact with fine branching, making it ideal for high\--detail bonsai.

    \item \textbf{Hardiness:} Hardy to around \--10\textdegree C, but benefits from some winter protection in temperate climates.
\end{itemize}

Its suitability for fine deadwood carving and jin/shari work has made it a favourite among bonsai artists worldwide.

\medskip
\topic{Why Choose Itoigawa? Advantages and Considerations}
\noindent
\textit{Itoigawa} has earned its popularity due to:

\begin{itemize}
    \item \textbf{Vigour:} Responds well to pruning and backbuds reliably on old wood.

    \item \textbf{Appearance:} Its foliage pads can be refined to a very high degree without appearing coarse or overgrown.

    \item \textbf{Size Versatility:} Suitable for all sizes, from shohin to large bonsai.

    \item \textbf{Styling Flexibility:} Adapts well to formal and informal upright, semi\--cascade, and literati styles.
\end{itemize}

That said, it requires careful watering and good airflow to avoid fungal issues, and styling should respect its natural foliage density.

\medskip
\topic{Soil and Potting Requirements}
\noindent
As with most junipers, free\--draining soil is essential.
A standard mix might include:
\begin{itemize}
    \item \textbf{60\% Akadama:} for water retention and cation exchange.

    \item \textbf{20\% Pumice:} for root oxygenation.

    \item \textbf{20\% Lava Rock or Kiryu:} for structure and long\--term aeration.
\end{itemize}

Slip\--potting (without disturbing roots) is safe almost any time if done with care.
Full repotting should be reserved for early spring, and root pruning should be cautious and spaced over several years.

\medskip
\topic{Light and Temperature Preferences}
\noindent
\textit{Itoigawa} thrives with full sun exposure but may benefit from afternoon shade in extreme summer heat.
Proper lighting is essential for maintaining compact growth and foliage colour.

\begin{itemize}
    \item \textbf{Sunlight:} 6+ hours of direct light daily.

    \item \textbf{Wintering:} Hardy, but shelter from wind and severe frost (\textless{} \--5\textdegree C) is recommended.

    \item \textbf{Ventilation:} Essential year\--round to prevent fungal problems and spider mites.
\end{itemize}

\medskip
\topic{Watering and Humidity}
\noindent
Junipers dislike sitting in water.
Overwatering leads to root rot, yet underwatering causes foliage dieback, which often appears much later than the stress event.

\begin{itemize}
    \item \textbf{Check daily:} Water when the top layer of soil begins to dry.

    \item \textbf{Water deeply:} Until water flows freely from drainage holes.

    \item \textbf{Humidity:} Beneficial but not essential. Misting can help in hot dry spells.
\end{itemize}

Avoid watering late in the evening during humid months---fungal pathogens thrive in still, wet conditions.

\medskip
\topic{Fertilisation Strategy}
\noindent
\textit{Itoigawa} appreciates a balanced feeding regimen. Organic fertilisers like rapeseed pellets are well tolerated.

\begin{itemize}
    \item \textbf{Spring through early summer:} Use balanced (e.g., 10\--10\--10) or nitrogen\--heavy fertiliser.

    \item \textbf{Late summer to autumn:} Switch to a low\--nitrogen blend (e.g., 4\--10\--10) to aid lignification.

    \item \textbf{Winter:} Avoid fertilising; tree is semi\--dormant.
\end{itemize}

Weak trees or freshly styled specimens should be fed cautiously to avoid stressing them.

\medskip
\topic{Propagation Techniques}
\noindent
\textit{Itoigawa} is typically propagated via cuttings, though grafting is used for high\--end stock.

\begin{itemize}
    \item \textbf{Semi\--hardwood cuttings:} Taken in late summer, treated with rooting hormone, and placed in coarse, sterile substrate.

    \item \textbf{Grafting:} Often used to add Itoigawa foliage to a stronger rootstock or improve nebari.
\end{itemize}

Air\--layering is rarely used on junipers, as they don't reliably produce roots through cambial removal.
Patience and attention to hygiene are key.

\medskip
\topic{Pruning and Styling Techniques}
\noindent
Styling \textit{Itoigawa} is a joy but requires finesse:

\begin{itemize}
    \item \textbf{Structural pruning:} Best done in late winter or early spring.

    \item \textbf{Pinching:} Remove new tips gently between fingernails; avoid scissors to prevent browning.

    \item \textbf{Wiring:} Can be done anytime, but the tree should not be under stress.
    Check regularly to prevent wire bite.

    \item \textbf{Jin and shari:} Itoigawa tolerates deadwood well; use lime sulfur to preserve whitened features.
\end{itemize}

\medskip
\topic{Training Timeline Examples}
\noindent
\textbf{Year 1:} Acquisition and slip\--potting into a well\--drained training pot.
Light wiring if needed.
Allow unrestricted growth.

\textbf{Year 2:} Select primary branches and do first major wiring and jin work.
Fertilise generously.
Begin trunk shaping with guy wires.

\textbf{Year 3:} Pinch to develop pads.
Begin reducing strong leaders to divert energy.
Establish foliage silhouette.

\textbf{Year 4+:} Annual refinement.
Continue pad formation and detail wiring.
Begin transitioning to final bonsai pot.

\medskip
\topic{Comparative Cultivar Notes}
\noindent
How does \textit{Itoigawa} compare to other junipers?

\begin{itemize}
    \item \textbf{vs. Kishu:} Kishu has denser, slightly more waxy foliage and is slower growing.
    More compact but less dramatic deadwood potential.

    \item \textbf{vs. Shimpaku:} Classic Shimpaku has slightly looser foliage and is more vigorous but less visually refined.

    \item \textbf{vs. San Jose:} San Jose juniper reverts to needle foliage more easily and has a coarser overall look.
\end{itemize}

\medskip
\topic{Common Mistakes to Avoid}
\noindent
\begin{itemize}
    \item \textbf{Overwatering:} Always check soil moisture first.

    \item \textbf{Under\--fertilising in early years:} Trees need energy to build structure.

    \item \textbf{Pinching too early:} Allow growth to elongate slightly before pinching to maintain vigour.

    \item \textbf{Overworking:} Avoid combining repotting and styling.
    Give ample rest.

    \item \textbf{Neglecting airflow:} Leads to fungal issues, especially in humid climates.
\end{itemize}

\medskip
\topic{In Closing}
\noindent
Whether you're planning to build a future showstopper or simply want to hone your refinement techniques, \textit{Itoigawa} is a responsive, beautiful, and rewarding bonsai species.
The cultivar's ability to take detailed styling work, combined with its elegant foliage, has made it a top choice for professionals and hobbyists alike.

\textit{Have questions or additional techniques to share? Members can suggest improvements or request monographs on other species by joining the club WhatsApp group.}
\closearticle
\clearpage
